% 
%% SOSP 2017 Template
%%
%% Uses sigplanconf from:
%% 
%%    http://www.sigplan.org/sites/default/files/sigplanconf.cls
%%
%% with 10pt and preprint options. 
%%
%% Replace 'XX' with your paper number (assigned when you register abstract)
%% Replace 'NN' with actual number of pages. 

\documentclass[10pt,preprint,nocopyrightspace]{sigplanconf}
\usepackage{times}

\usepackage{datetime}
\usepackage{url}
\usepackage{hyperref}
\usepackage{graphicx}
\usepackage{grffile}
\usepackage[T1]{fontenc}
\usepackage{listings}
\usepackage{float}
\usepackage{natbib}


\bibliographystyle{IEEEtran}
\setcitestyle{numbers,square,comma, sort&compress}


% These only appear when the 'preprint' option is specified.
% Enabling these will cause the first page of the document to fail the 
% format check on HotCRP :-(
%\titlebanner{Under submission to SOSP 2017 - do not cite or distribute}
%\preprintfooter{Draft of {\currenttime}, \today{}}

% No date in title area.
\date{}

% Paper number and no. of pages as author
\authorinfo{\textbf{Shao-Jung (Louie) Lu, Ching-Chun (Jim) Huang}}


% Actual document begins below.
\begin{document}

\title{F9: A secure microkernel for MMUless embedded systems}
\maketitle

% \tableofcontents


\begin{abstract}\textsl{}
	
Internet of Things (IoT) can be seen in such as vehicles, sensors nodes, smartphones, and medical appliances nowadays. The hardware of IoT is typically constrained regarding computing power, available memory, communication bandwidth, and energy capacities, which implies the whole system should be expected to fulfill these requirements: reliability, real-time behavior.

Microkernels are suitable for IoT constraints in concept, but its approaches such as seL4 and Fiasco.OC, still need a rich hardware resource with MMU for memory management and higher power consumption.

In this paper, we describe the design and implementation of F9 microkernel - a microkernel based on L4, a family of second-generation microkernels, intended to build for deeply embedded IoT devices for efficiency, security, and flexible development. F9 microkernel is optimized for ARM Cortex-M series microprocessors, with only 3500 lines of source code, which is the smallest L4-compatible microkernel as far as we know.

\end{abstract}

\section{Introduction}

Nowadays IoT devices are ubiquitous in our life, most of them are connected to the Internet. In this context, security and reliability concerns become important issues since IoT applications today often deal with sensitive data and act on the physical world.

Research in operating systems shows that monolithic kernel such as Linux kernel\cite{chou2001empirical} and Microsoft Windows NT\cite{swift2005improving} are making errors from drivers, and Gernot Heiser, who was doing research in microkernel, state that larger trusted computing base (TCB) cannot be made trustworthy, the key to a trustworthy TCB is a very small kernel. Compare to the monolithic kernel, microkernel is more suitable for IoT devices since microkernel will only contain code that must be privileged\cite{heiser2005secure}.

In microkernel, kernel only provides three mechanisms: IPC, Memory Management and thread management, and these mechanisms can be used to build up the services, for example, file systems or device drivers, which can provide complete functions for unprivileged applications\ref{fig:kernel_diff}.

The history of microkernel can be date back to 1969s, and have developed in 3 generations: (i) Mach, (ii) L4, improving IPC performance, originally written in assembly (iii) seL4, Fiasco.OC, NOVA, platform independence and formal verification, security, SMP, etc. These microkernels need to run on the platform which has Memory Management Unit, but in low-power consumption IoT device, it lacks MMU.

The ARM Cortex-M series microprocessors are the most widely used in low-power consumption IoT devices, as for 2015, it has shipped more than 10 billion Cortex-M in the market\cite{arm2015cortex}. These processors has high performance and low-power consumption but were lack of conventional Memory Management Unit, which limits their ability to run traditional operating systems with (paged) memory protection. Instead, F9 microkernel is designed to use the Cortex-M Memory Protection Unit to provide memory isolation. 

\begin{figure}[H]
	\begin{center}
		\includegraphics[width=\linewidth]{picture/kernel_diff.png}
	\end{center}
	\caption{kernel space service is split out to user space in microkernel \cite{heiser2012}}
	\label{fig:kernel_diff}
\end{figure}


\section{Background of microkernel}

\subsection{Brief history of microkernel}

The first generation of microkernel "Mach", demonstrate that the idea of Inter-process communication (IPC) was successful to increase kernel modularity, security, and scalability, it replace pipes with IPC, and split other services out from kernel to improved stability of microkernel, although IPC performance is vital for microkernel, however, limited to the implementation of Mach, the performance of IPC was poor in modern processor\cite{liedtke1993improving}.

The second generation of microkernel "L3 \& L4 family", start from Liedtke\cite{liedtke1993improving}, L3 kernel tries to improve IPC performance by re-design microkernel from architecture to coding level, the result of L3 kernel improve IPC performance nearly twentyfold, Liedtke then introduce the "L4" with "V2" ABI in the microkernel community from 1995\cite{liedtke1995u}. This work triggered a twenty-year evolution, "L4" then evolved into a family, including the first ported L4 kernel \textit{Hazelnut} in 2000 (from Pentium to ARM), \textit{L4Ka::Pistachio} that focus on portability, Pistachio has ported into MIPS, Alpha, 64-bit PowerPC and ARM, and \textit{Fiasco} that aim for formal verification.

\subsection{Inter-Process Communication}

Inter-Process Communication (IPC) is the mechanism that microkernel thread communicate with other threads\cite{liedtke1993improving}. IPC is a critical issue in microkernel since operating system services such as file systems, device drivers, network protocol stacks, etc, are not include in the kernel but running in the user space, so, all the communication between system services and user applications involves IPC to and from servies\cite{iqbal2010overview}, typically message passing. IPC may design in synchronous or asynchronous, but in the principle of minimality, a kernel will only choose one approach. The L4 kernel provides message-based, synchronous IPC between threads for performance reason since Liedtke research in 1993\cite{liedtke1993improving}, only OKL4 used the asynchronous IPC approach\cite{elphinstone2013l3}.

\subsection{Memory Management and Address spaces}

In modern system architectures and operating system, physical memory usually will be 'virtualize' into virtual memory for programs to execute in\cite{arpaci2015operating}. A hardware device called memory-management unit (MMU) will translate between physical memory address to virtual memory and so on are lies between CPU and memory.

L4 provides the mechanism for implementing recursively-defined virtual address that is managed in user level\cite{dannowski2011l4}. Initial address space is controlled by the first process, other processes which derive from the first process will obtain memory pages from first.

\subsection{Thread Management}

Threads are the basic execution unit in L4. They are created, destroyed via the \verb|Thread-Control| system call\cite{nourai2005aphysically}. The thread dispatcher is responsible for switching contexts. Each thread has its own Thread Control Block (TCB) and addressed by its global id, and the TCB is where a thread's state is saved to or restored from on a context switch\cite{nourai2005aphysically}.

\subsection{Thread identifiers}

Every threads in L4 have \textit{global} and \textit{local} thread identifier (thread id). Global thread id is unique throughout the entire system, and local thread id has the scope limited in its own address space.

The global thread identifier in L4 contains two distinct parts - an upper 18-bits thread number and lower 14-bits version number. Also, there are two special thread ids reserved L4 API - NilThread and AnyThread. NilThread identifier is implemented as 0, and AnyThread is implemented as all bits set to 1. These special thread identifier will be used in L4 IPC system call\cite{nourai2005aphysically}.

\section{Design of F9 microkernel}

The F9 microkernel aims to run on the deeply embedded devices, and its goals are to make itself efficient and secure.

\begin{figure}[H]
	\begin{center}
		\includegraphics[width=\linewidth]{picture/f9_arch.png}
	\end{center}
	\caption{F9 microkernel system architecture}
	\label{fig:f9}
\end{figure}

\subsection{Memory Management}

Unlike other L4 kernels, F9 microkernel focused on the MPU-only platform such as ARM Cortex-M, especially optimized for Cortex M3/M4, ARM Cortex-M supports the ARMv7 Protected Memory System Architecture (PMSAv7) model. The system address space of a PMSAv7 compliant system is protected by a Memory Protection Unit(MPU). Also, the RAM was small (256 Kbytes) but a larger Physical address space (up to 32-bit) was used, including the hardware device, flash, and bit-band zone\cite{yiu2013definitive,st2016managing,usna2015l18}.

The MPU protected memory is divided up into a set of regions, with the number of regions supported IMPLEMENTATION DEFINED, for example, in STM32F429, it provides eight separate memory regions. In PMSAv7, the minimum protect region size is 32 bytes, and maximum up to 4 GB. The MPU provides full support for:
\begin{itemize}
	\item Protection regions
	\item Overlapping protection regions
	\item Access permissions
	\item Exporting memory attributes to the system
\end{itemize}

MPU mismatches and permission violations invoke the programmable priority MemManage fault handler.

There is a limitation of setting memory region in MPU is that developer can only set "$2^n$" size memory region from 32 bytes to 4 Gbytes, and must be aligned by the size\cite{yiu2013definitive,st2016managing}. For example, setting 64 bytes memory region is valid, but setting 96 bytes memory region is invalid, and a 64-KB sized region must be aligned on a multiple of 64KB, for example, \verb|0x000100000| or \verb|0x000200000|. Obey to this limitation, F9 microkernel developed an elegant mechanism called fpage, which the size it fit to set in MPU regions.

Memory management in F9 microkernel, can split into three conceptions: 1. memory pool, which represents the area of PAS with specific attributes (hardcoded in mem map table). 2. address space - sorted list of fpages bound to particular thread(s). 3. flexible page - unlike traditional pages in L4, fpage represent in MPU region instead. 

\subsubsection{Memory Pool}
In F9 microkernel, memory pool was static and declare before compile; it will declare each memory segment's start, end, permission, etc., then slice the needed part into fpage size, put into address space.

\subsubsection{Address Space}
Address space is the structure that each thread record their fpage. When creating a new thread in F9 microkernel, it will first declare fpage chain for it, then mapping needed information into address space.

In F9 microkernel, address space is using the recursive address space model from L4, and allows to construct and manipulate address spaces with the following three operations: \verb|grant|, \verb|map|, and \verb|unmap|\cite{volp2002design}.

Kernel can used \verb|map_area| function to manipluate address spaces, \verb|grant| operation means a fpage p is mapped from address space $\sigma{a}$ to $\sigma{b}$, afterwards, p is only accessible in $\sigma{b}$ address space\ref{fig:grant}.

\begin{figure}[H]
	\begin{center}
		\includegraphics[height=1.5in]{picture/grant.png}
	\end{center}
	\caption{granting fpage p from $\sigma{a}$ to $\sigma{b}$.}
	\label{fig:grant}
\end{figure}


\verb|map| operation meas a fpage p is mapped from address space, afterwards, p is accessible in both address spaces in $\sigma{a}$ and $\sigma{b}$, \ref{fig:map}.

\begin{figure}[H]
	\begin{center}
		\includegraphics[height=1.5in]{picture/map.png}
	\end{center}
	\caption{mapping fpage p from $\sigma{a}$ to $\sigma{b}$.}
	\label{fig:map}
\end{figure}


\verb|unmap| operation will unmapped from all address space directly or indirectly recieved it from $\sigma{a}$, afterwards, p is no longer accessible in $\sigma{b}$, $\sigma{c}$ and $\sigma{d}$\ref{fig:unmap}.

\begin{figure}[H]
	\begin{center}
		\includegraphics[height=1.5in]{picture/unmap.png}
	\end{center}
	\caption{unmapping fpage p in $\sigma{a}$, $\sigma{b}$, $\sigma{c}$ and $\sigma{d}$ were no longer accessible to fpage p.}
	\label{fig:unmap}
\end{figure}



\verb|map_area| can find the proper address space from source to destination by giving base position and size, it will search the first and last fpage that contain base position and end of the request size. If the first and last fpage were found, it will split fpage into two fpage, then mapping fpage from source to destination. Also, F9 microkernel provides these API for low-level address space settings:
\begin{itemize}
\item \verb|as_create|
\item \verb|as_setup_mpu|
\end{itemize}

\verb|as_create| is used when creating a new thread, and it will allocate a new address space from address space table.

\verb|as_setup_mpu| is used when to perform a context switch between two threads.


\subsubsection{Flexible Pages (fpage)}
In F9 microkernel, every fpage can directly use by MPU as a page (that mean be protected by MPU), and address space was built from fpage. Fpage structure includes base address, memory pool id, flags, size and permission. F9 microkernel provides fpage API:
\begin{itemize}
	\item \verb|create_fpage|
	\item \verb|create_fpage_chain|
	\item \verb|assign_fpages|
\end{itemize}

\subsection{IPC}

L4 provides a single system call for all IPC functions, F9 microkernel follows the implementation from L4\cite{dannowski2011l4}. The IPC was implemented in synchronous IPC, means that when a thread wants to communicate with another thread, it will first set sender into the \verb|SEND_BLOCKED| state, then waiting for target receiver goes into the \verb|RECV_BLOCKED| state. After these two threads are in the proper state, it will thus perform IPC communication. Synchronous IPC avoids buffering in the kernel, the management and the copying cost between two thread.\cite{nourai2005aphysically}

IPC are not only providing threads communication, but also handle the protocol for user thread create, user interrupt config and handle, and user space standard output. The method that IPC can distinguish different function is to check \verb|from_tid| and \verb|to_tid|, as mentioned earlier, there was two special thread tid in L4 kernel  called \verb|NilThread| and \verb|AnyThread|, F9 microkernel also provide other special thread tid to perform different functions.

F9 microkernel introduces two different IPC register - message register and buffer register. The message register is a virtual register that can only be read one time, after one read, the value in this message register will be undefined. First introduced the concept in pistachio\cite{nourai2005aphysically}, Message register can be moderate-size, and now F9 microkernel implements 4 type of items: Untyped, Typed, Map, Grant item.

\subsection{Threads}

Threads are the basic execution unit in F9 microkernel, each thread in F9 microkernel, has its own Thread Control Block (TCB) and addressed by its global id. In F9 microkernel, global id only used the upper 18-bits and lower bits are reserved.

In F9 microkernel, there is a thread map which stores the TCBs, and sorted by \textit{global id} in ascending order. When the thread was created, it will insert into the thread map by its global id, the benefit of ensuring the thread map is in ascending order is that F9 microkernel can get the thread by global id using binary search, provide the API \verb|thread_by_globalid| for other kernel components.

\subsection{Scheduler}

The thread scheduler has the responsibility to select new thread for context switches when the current thread was blocked or exit. In F9 microkernel, scheduler perform as a round-robin scheduler with non-priority, and tries to select the first runnable normal thread, for minimality, F9 microkernel didn't implement runnable timeslice for threads. This approach has a drawback that if user space thread was in a busy loop, it will hang on busy loop thread and never transfer to other normal thread.

\begin{figure}[H]
	\begin{center}
		\includegraphics[width=\linewidth]{picture/scheduler.png}
	\end{center}
	\caption{F9 microkernel scheduler hierarchy}
\end{figure}

The future work on F9 microkernel thread and the scheduler will need to implement thread timeslice for preventing thread scheduling starvation and to search the real-time scheduling algorithm.


\subsection{Ksymbol maps and Kprobe}
F9 microkernel using an elegant method to dump out the symbol from kernel and let kernel can using these symbol information such as function address and symbol name in runtime. When building the entire kernel, the build system will first compile all component then linked to a non-symbol elf kernel file, then the special make scripts which build to covert elf to symbol map file, will using GNU toolchain "nm" tool to dump all symbols out of the kernel into a pre-define c template. After all, it compiled and linked the non-symbol kernel and symbol map which set into the \texttt{.sym\_tab} section.

This mechanism lets F9 microkernel can easily using symbol maps in runtime, and the demonstration about using symbol map can be found at Kprobe.

F9 microkernel has implemented a light weight debugging system call "Kprobe", this term is inherited from Linux kernel. The Kprobe goals are to provide an in-kernel dynamic tracing mechanism, allow developers to gather additional information about kernel operation without recompiling or rebooting system. Currently, Kprobe is implemented through hardware breakpoint from ARMv7 Debug Architecture. 

In ARMv7-M, we can use DebugMon exception to trap probed address\cite{yiu2013definitive}, when a kprobe is registered, one of the FPB (Flash patch and Breakpoint Unit) will set to the probed address. When CPU hits the probed address, the trap occurs, FPB will generate a DebugMon exception and set bit-BKPT in Debug Fault Status Register (DFSR).

After all \texttt{pre\_handlers} of kprobes associated with the probed address are called, Kprobes sets \texttt{bit-MON\_STEP} in Debug Exception and Monitor Control Register in order to single-step the probed address; Then Kprobes returns to
the probed address. And soon a DebugMon exception will be generated again once an instruction is executed.

We can check bit-HALTED from DFSR to know if there is a single-step debug event. If yes, it is time to call all the \texttt{post\_handlers} associated with this address. Then execution resumes and continues normally.

Because of the hardware support, we can use 10 times less code compare to Linux kernel to implement Kprobes. But, at this moment, F9 microkernel Kprobe can only be used in kernel space, user space was not supported yet, cause F9 microkernel is using function address to regist in Kprobe, compare to Linux kernel, it is using symbol name to register the probed function. But as mention below, F9 microkernel already implement the infrastructure to support symbol name register the probed function in user space, this will be the next develop in F9 microkernel.

\begin{lstlisting}[basicstyle=\small,frame=single]
/* Linux kernel */
struct kprobe kp = {
	.symbol_name    = "_do_fork",
};

/* F9 microkernel */
struct kprobe kp = {
	.addr    = ktimer_handler
};

\end{lstlisting}

\subsection{Ktimer}
F9 microkernel implement Ktimer for handle timer ticks interrupt, is provide an will interface for kernel to create ktimer event, kernel will set the ticks and event function which it wants to execute after how many ticks, when ticks were down to 0, ktimer will trigger an ktimer softirq to execute function handler in ktimer event.

\subsection{Tickless scheduling}

F9 microkernel also provides tickless scheduling to prevent higher power consumption in deeply embedded system\cite{freertos1,freertos2}. In Cortex-M4, it has a 24-bit system timer, SysTick, that countdown from the reload value to zero then trigger timer interrupt, then reload the value in the \verb|STK_LOAD| register on the next clock edge\cite{st2016manual}. Assume we config \verb|STK_LOADS| value to 1024, then the SysTick will generate Timer interrupt every 1024 ticks.

In tickless scheduling mode, kernel will enter tickless right before doing to CPU idle state, set interval of next timer interrupt as a delta of next event, or \verb|KTIMER_MAXTICKS|. Then adjust system time after waked up.

\begin{figure}[H]
	\begin{center}
		\includegraphics[width=\linewidth]{picture/tickless.png}
	\end{center}
	\caption{Tickless scheduling in F9 microkernel}
\end{figure}

\section{Implementation techniques}

While implementing microkernel for IoT devices, we will encounter hardware constraint such as limited RAM space, slower CPU clock rate, etc, we will need to use some techniques for microscope optimization.

\subsection{sorting}
F9 microkernel sort routine was taken from BSD\cite{bentley1993engineering} developed in University of California, Berkeley, BSD approach found that swapping is the most sensitive issue in sorting. It have conquered the problem that sorting in different width data size. Also, F9 microkernel implement sort by heapsort to prevent quick sort worst case in $\mathcal{O}(n \times n)$ , make the sorting time is $\mathcal{O}(n\log{}n)$  both on average and worst-case.

\subsection{memcpy, duff's device}
F9 microkernel memcpy was derive from MUSL libc project\cite{felker}, it using the technique from "Duff's device", also called "loop unrolling" to improve the performance of memcpy. The future work is to rewrite memcpy code in ARM Thumb-2 to improve performance.

\subsection{Init hooks}

F9 microkernel used a technique call "Initialize hooks", this technique can let developer write code in kernel and assign to run the code in the specific stage when kernel is initializing.

API \verb|INIT_HOOK(hook, level)| is the syntactic sugar that will convert into a const \verb|init_struct|, which will then put in section \verb|.init_hook|, in linker script it will define \verb|init_hook| start and end position, it will be a linear address for all init hook defined by developer.

In booting stage, kernel will try to execute these hook by \verb|run_init_hook(level)|, this function will then scan the \verb|init_hook| region from start to end, if it found the proper level to execute, it will then execute the hook. Also, when multiple hooks are on the same level, this approach won't guarantee the order to execute hooks.

\subsection{bit-band and Ktable}
ARM Cortex-M provide an atomic access method call "Bit-banding", Bit-banding mechanisms is the device takes a region of memory (the Bit-band region) and maps each bit in that region to an entire word in a second memory region (the Bit-band Alias Region).\cite{schaenzle2013}

\begin{figure}[H]
\begin{center}
	\includegraphics[width=\linewidth]{picture/bit-banding.png}
\end{center}
\caption{bit-banding in Cortex-M\cite{st2016manual}}
\end{figure}

This technique makes atomic access more friendly to developer since developer does not need to write in assembly nor take care about context-switch or preemption.

F9 microkernel using bit-band technique to create a fast atomic object management mechanism "Ktable", Ktable can create an structure that provides table size, number of chunks, and data. F9 microkernel using Ktable to manage fpages table, address space table, and thread table.

\begin{itemize}
\footnotesize
\item \verb|void ktable_init(ktable_t *kt);|
\item \verb|int ktable_is_allocated(ktable_t *kt, int i);|
\item \verb|void *ktable_alloc_id(ktable_t *kt, int i);|
\item \verb|void *ktable_alloc(ktable_t *kt);|
\item \verb|void  ktable_free(ktable_t *kt, void *element);|
\item \verb|uint32_t ktable_getid(ktable_t *kt,void *ele);|
\end{itemize}


\begin{figure}[H]
	\begin{center}
		\includegraphics[width=\linewidth]{picture/ktable.png}
	\end{center}
	\caption{Manipulation in Ktable\cite{ncku2015}}
\end{figure}

\subsection{Mutual exclusion lock}
ARM Cortex-M provide a local exclusive monitor for exclusive accesses. This gave Cortex-M the atomic load / store ability for developers\cite{arm2012v7}. F9 microkernel implements a POSIX compatible library for user space, using atomic load / store technique, we are able to implement a userspace mutex lock implement via POSIX pthread interface.

Using thumb instruction, our approach in F9 microkernel only need 24 bytes to implement \texttt{pthread\_mutex\_trylock} function in pthread library.

\begin{lstlisting}[basicstyle=\small,frame=single]
44:	4603      	mov	r3, r0
46:	f04f 0101 	mov.w	r1, #1
4a:	461a      	mov	r2, r3
4c:	e852 0f00 	ldrex	r0, [r2]
50:	2800      	cmp	r0, #0
52:	bf04      	itt	eq
54:	e842 1000 	strexeq	r0, r1, [r2]
58:	4603      	moveq	r3, r0
5a:	4618      	mov	r0, r3
5c:	4770      	bx	lr
\end{lstlisting}

\section{Conclusion}
F9 microkernel is the first microkernel that using MPU and create security memory management microkernel, rely on the fpage and address space infrastructure, F9 microkernel only wastes a little memory space but provide the thread to thread memory access control. Also, F9 microkernel have a fast IPC performance and tickless schedule for energy-saving and provide F9 microkernel a well in-kernel debugging toolchains like kdb and Kprobe.

The future work of F9 microkernel can split into two categories. One is to develop kernel itself, and another is to improve kernel dynamic tracing toolchains.

\bibliography{f9-report.bib}

\end{document}
